%%%%%%%%%%%%%%%%%%%%%%%%%%%%%%%%%%%%%%%%%%%%%%%%%%%%%%%%%%%%%%%%%%%%%%%%%%%
% Función: cálculo de la calificación final de los formularios.
%
% Uso: mantenimiento de la plantilla; normalmente no debes modificarlo.
%%%%%%%%%%%%%%%%%%%%%%%%%%%%%%%%%%%%%%%%%%%%%%%%%%%%%%%%%%%%%%%%%%%%%%%%%%%

% Las notas de entrada son valores sobre diez, no aportaciones ponderadas.
\newcommand{\CheckPaperworkGrade}[2]{%
  \pgfmathtruncatemacro{\gradeOutsideRange}{(#1 < 0) || (#1 > 10)}%
  \ifnum\gradeOutsideRange=1
    \PackageError{grades-calculation}{#2 must be between 0 and 10}%
      {Use a numerical mark out of ten in Config/myconfig.tex.}%
  \fi}

% La aritmética decimal evita errores de clasificación en límites como 7 y 9.
\ExplSyntaxOn
\newcommand{\CalculateTFGFinalGrade}{%
  \edef\myFinalGrade{\fp_eval:n {(7 * \varA + 3 * \varB) / 10}}}
\ExplSyntaxOff

\pgfmathsetmacro{\varB}{\myTribunalGrade}
\CheckPaperworkGrade{\varB}{myTribunalGrade}
\ifthenelse{\equal{\myWorkType}{TFG}}
{%
  \pgfmathsetmacro{\varA}{\myAcademicTutorGrade}
  \CheckPaperworkGrade{\varA}{myAcademicTutorGrade}
  \CalculateTFGFinalGrade
  \pgfmathtruncatemacro{\gradeContradictsFailedComponent}{%
    ((\varA < 5) || (\varB < 5)) && (\myFinalGrade >= 5)}
  \ifnum\gradeContradictsFailedComponent=1
    \PackageError{grades-calculation}{TFG weighted pass with a failed component}%
      {Article 9.3 requires a failed TFG when either component is below five.^^J%
       Ask the department how to record the final mark; no capped mark is invented.}%
  \fi
}
{%
  \ifthenelse{\equal{\myWorkType}{TFM}}
  {\edef\myFinalGrade{\varB}}
  {\PackageError{grades-calculation}{Unsupported work type: \myWorkType}%
     {This helper is intended for TFG and TFM paperwork only.}%
   \def\myFinalGrade{0}}
}

\pgfmathparse{
  int( (\myFinalGrade < 5) * 1 +
  (\myFinalGrade >= 5) * (\myFinalGrade < 7) * 2 +
  (\myFinalGrade >= 7) * (\myFinalGrade < 9) * 3 +
  (\myFinalGrade >= 9) * 4 )
}
\edef\gradeCategory{\pgfmathresult} % Store as an integer

% Convert numeric category to text
\def\myFinalGradeName{%
  \ifnum\gradeCategory=1 suspenso%
  \else\ifnum\gradeCategory=2 aprobado%
  \else\ifnum\gradeCategory=3 notable%
  \else sobresaliente%
  \fi\fi\fi
}


\ifthenelse{\equal{\myTribunalMHProposal}{YES}}
  {
    \newcommand{\MHProposalBoxes}{$\XBox$ Sí ~~~~ $\Box$ No}
  }
  {
    \newcommand{\MHProposalBoxes}{$\Box$ Sí ~~~~ $\XBox$ No}
  }
  
